\section{The shared-factor character gate}
\label{sec:tpc21-character}

The additive transform in \eqref{eq:tpc21-long-packet} is naturally
defined on the full cyclic group modulo \(q\).  The coefficient itself,
however, is supported on the unit group.  Multiplicative Fourier
analysis on that unit group retains the provenance of both rows and,
in particular, keeps the generic Schur mask inside the transform.  This
section gives the exact change of basis.  It is an identity, not an
estimate for the resulting character moment.

For a positive integer \(s\), write

\[
 G_s=(\mathbb Z/s\mathbb Z)^\times,
 \qquad
 \mathfrak X(s)=\widehat G_s,
\]

with \(G_1\) and \(\mathfrak X(1)\) interpreted as trivial groups.  A
character is extended by zero away from the units whenever it is
evaluated on an integer.  For \(f:G_s\to\mathbb C\), use the convention
\begin{equation}
 \widehat f^{\,\times}(\chi)
 =\sum_{k\in G_s}f(k)\overline{\chi(k)}.
 \label{eq:tpc21-multiplicative-transform}
\end{equation}
Then
\begin{align}
 f(k)
 &=\frac1{\varphi(s)}
   \sum_{\chi\in\mathfrak X(s)}
   \widehat f^{\,\times}(\chi)\chi(k),
 \label{eq:tpc21-multiplicative-inversion}\\
 \sum_{k\in G_s}|f(k)|^2
 &=\frac1{\varphi(s)}
   \sum_{\chi\in\mathfrak X(s)}
   |\widehat f^{\,\times}(\chi)|^2.
 \label{eq:tpc21-multiplicative-parseval}
\end{align}
In particular, if
\(f^\circ=f-|G_s|^{-1}\sum_{k\in G_s}f(k)\), then
\begin{equation}
 \boxed{
 \|f^\circ\|_{2,G_s}^2
 =\frac1{\varphi(s)}
  \sum_{\substack{\chi\in\mathfrak X(s)\\\chi\ne\mathbf1_s}}
  |\widehat f^{\,\times}(\chi)|^2.}
 \label{eq:tpc21-nonprincipal-parseval}
\end{equation}
Thus ``connected'' below means that the principal unit-group character,
equivalently the fiber mean, has been removed.

\subsection{The masked shared-\texorpdfstring{\(g\)}{g} convolution}

Fix squarefree divisor moduli \(u,v\leq R\) with \((uv,H)=1\).
Modulus pairs failing this coprimality condition have empty admissible
fibers and are omitted.  Compatibility of the two row labels is imposed
inside the fiber.  Write
\begin{equation}
 g=(u,v),\qquad u=ga,\qquad v=gb,\qquad q=[u,v]=gab.
 \label{eq:tpc21-character-gab}
\end{equation}
The integers \(g,a,b\) are pairwise coprime.  CRT therefore identifies
\(G_q=G_g\times G_a\times G_b\).  Write a character of \(G_q\) as

\[
 \chi=\chi_g\boxtimes\chi_a\boxtimes\chi_b,
 \qquad
 \chi_c\in\mathfrak X(c).
\]

For characters \(\alpha\in\mathfrak X(u)\) and
\(\beta\in\mathfrak X(v)\), define the masked two-row
Dirichlet-character sum
\begin{equation}
 \mathcal B_{\mathfrak m}^{u,v}(\alpha,\beta)
 =\sum_{\omega_1,\omega_2\in\mathcal A}
 x_{\omega_1}y_{\omega_2}
 \mathfrak m(\omega_1,\omega_2)
 \alpha(m_{\omega_1})\beta(m_{\omega_2}).
 \label{eq:tpc21-masked-B}
\end{equation}
Here \(\omega_i=(\ell_i,d_i)\) is a row label.  The zero extension of
the characters inserts the row--modulus coprimality conditions.  There
is no compatibility indicator in \eqref{eq:tpc21-masked-B}; it will be
created exactly by averaging a character of the shared factor \(g\).

\begin{theorem}[Exact masked shared-factor transform]
\label{thm:tpc21-masked-shared-character}
With the conventions above,
\begin{equation}
 \boxed{
 \widehat Z_{u,v}^{\,\mathfrak m,\times}(\chi)
 =\frac{\overline{\chi(-h)}}{\varphi(g)}
  \sum_{\psi\in\mathfrak X(g)}
  \mathcal B_{\mathfrak m}^{u,v}
  \left(
   (\chi_g\psi)\boxtimes\chi_a,
   \overline\psi\boxtimes\chi_b
  \right).}
 \label{eq:tpc21-masked-shared-character}
\end{equation}
The factor \(1/\varphi(g)\) is essential.  It is the normalized
orthogonality projector onto compatible row pairs.
\end{theorem}

\begin{proof}
Suppose first that a pair of rows contributes to
\(Z_{u,v}^{\mathfrak m}(\kappa)\).  In the CRT coordinates of
\eqref{eq:tpc21-character-gab},
\[
 \kappa_g\equiv-h\overline{m_1}
             \equiv-h\overline{m_2}\pmod g,
 \quad
 \kappa_a\equiv-h\overline{m_1}\pmod a,
 \quad
 \kappa_b\equiv-h\overline{m_2}\pmod b.
\]
Consequently
\begin{equation}
 \overline{\chi(\kappa)}
 =\overline{\chi(-h)}
  \chi_g(m_1)\chi_a(m_1)\chi_b(m_2).
 \label{eq:tpc21-row-character-phase}
\end{equation}
For unit rows, character orthogonality gives
\begin{equation}
 \one_{\{m_1\equiv m_2\ (g)\}}
 =\frac1{\varphi(g)}
  \sum_{\psi\in\mathfrak X(g)}
  \psi(m_1)\overline{\psi(m_2)}.
 \label{eq:tpc21-shared-character-projector}
\end{equation}
Insert \eqref{eq:tpc21-row-character-phase} and
\eqref{eq:tpc21-shared-character-projector} in the definition of the
multiplicative transform, then interchange the finite sums.  The two
row variables carry respectively
\((\chi_g\psi)\boxtimes\chi_a\) and
\(\overline\psi\boxtimes\chi_b\), proving
\eqref{eq:tpc21-masked-shared-character}.
\end{proof}

The unmasked identity has an additional factorization.  Define the
one-row Dirichlet polynomials
\begin{equation}
 \mathcal A_x^{(s)}(\rho)
 =\sum_{\omega\in\mathcal A}x_\omega\rho(m_\omega),
 \qquad
 \mathcal A_y^{(s)}(\rho)
 =\sum_{\omega\in\mathcal A}y_\omega\rho(m_\omega),
 \qquad \rho\in\mathfrak X(s).
 \label{eq:tpc21-row-Dirichlet-polynomials}
\end{equation}
When no ambiguity is possible, the superscript \(s\) is suppressed.
The relation \(k\equiv-h\overline m\pmod s\) also gives the useful
one-row formula
\begin{equation}
 \sum_{k\in G_s}A_s^x(k)\overline{\rho(k)}
 =\overline{\rho(-h)}\,\mathcal A_x^{(s)}(\rho),
 \label{eq:tpc21-row-fiber-character}
\end{equation}
and likewise for \(y\).  Indeed,
\(\overline{\rho(-h\overline m)}
=\overline{\rho(-h)}\rho(m)\).

\begin{corollary}[Unmasked row-polynomial factorization]
\label{cor:tpc21-unmasked-character-factorization}
If \(\mathfrak m\equiv1\), then
\begin{align}
 \mathcal B_1^{u,v}(\alpha,\beta)
 &=\mathcal A_x^{(u)}(\alpha)
   \mathcal A_y^{(v)}(\beta),
 \label{eq:tpc21-unmasked-B-factorization}\\
 \widehat Z_{u,v}^{\,1,\times}(\chi)
 &=\frac{\overline{\chi(-h)}}{\varphi(g)}
   \sum_{\psi\in\mathfrak X(g)}
   \mathcal A_x^{(u)}
    ((\chi_g\psi)\boxtimes\chi_a)
   \mathcal A_y^{(v)}
    (\overline\psi\boxtimes\chi_b).
 \label{eq:tpc21-unmasked-shared-character}
\end{align}
No conjugation of the \(y\)-polynomial is implicit.  In a Hermitian
packet it is already contained in the choice
\(y_\omega=\overline{x_\omega}\).
\end{corollary}

For the generic mask \eqref{eq:tpc21-generic-mask}, the factorization
\eqref{eq:tpc21-unmasked-B-factorization} is unavailable: the
determinant, same-source, and divisor-gcd conditions jointly depend on
both row labels.  Formula \eqref{eq:tpc21-masked-shared-character}
remains exact because it never removes that mask.

\subsection{Lcm aggregation and the exact fourth-moment gate}

Fix an admissible squarefree \(q\), so that \((q,H)=1\).  For every
ordered pair \(u,v\leq R\) with
\([u,v]=q\), use the decomposition \eqref{eq:tpc21-character-gab} and
restrict a given \(\chi\in\mathfrak X(q)\) to the corresponding CRT
factors.  Define
\begin{align}
 \mathcal T_q^{\mathfrak m}(\chi)
 &:=\sum_{\substack{u,v\leq R\\{}[u,v]=q}}
 \frac{\vartheta_R(u)\vartheta_R(v)}{\varphi((u,v))}
 \sum_{\psi\in\mathfrak X((u,v))}
 \mathcal B_{\mathfrak m}^{u,v}
 \left(
  (\chi_g\psi)\boxtimes\chi_a,
  \overline\psi\boxtimes\chi_b
 \right).
 \label{eq:tpc21-aggregated-character-coefficient}
\end{align}
Linearity and \cref{thm:tpc21-masked-shared-character} give
\begin{equation}
 \boxed{
 \widehat Z_q^{\,\mathfrak m,\times}(\chi)
 =\overline{\chi(-h)}\,
  \mathcal T_q^{\mathfrak m}(\chi).}
 \label{eq:tpc21-aggregated-character-transform}
\end{equation}
Since \((q,h)=1\), the phase on the right has modulus one.

\begin{theorem}[Exact masked connected character moment]
\label{thm:tpc21-exact-character-gate}
Define
\begin{equation}
 \mathfrak F_q^{\mathfrak m}
 =\frac1{\varphi(q)}
  \sum_{\substack{\chi\in\mathfrak X(q)\\\chi\ne\mathbf1_q}}
  |\mathcal T_q^{\mathfrak m}(\chi)|^2.
 \label{eq:tpc21-masked-fourth-moment}
\end{equation}
Then
\begin{equation}
 \boxed{
 \mathfrak F_q^{\mathfrak m}
 =\|D_q^{\mathfrak m}\|_{2,\mathcal U_q}^2.}
 \label{eq:tpc21-character-gate-identity}
\end{equation}
Thus the inherited sufficient condition
\eqref{eq:tpc21-inherited-gate} is exactly a bound for
\begin{equation}
 \sum_{q\in\mathcal Q}
 \bigl(\mathfrak F_q^{\mathfrak m}\bigr)^{1/2}.
 \label{eq:tpc21-family-character-gate}
\end{equation}
For example, with the epsilon bookkeeping of TPC-20, it suffices to
prove
\begin{equation}
 \sum_{q\in\mathcal Q}
 \bigl(\mathfrak F_q^{\mathfrak m}\bigr)^{1/2}
 \ll X^{-\eps}\frac{XQ}{\sqrt J}(\log X)^{-B}.
 \label{eq:tpc21-sufficient-character-gate}
\end{equation}
This is a sufficient input for the chosen short--short discrepancy
packet, not a necessary condition for cancellation of the scalar
correlation.
\end{theorem}

\begin{proof}
Apply \eqref{eq:tpc21-nonprincipal-parseval} to
\(f=Z_q^{\mathfrak m}\), and then use
\eqref{eq:tpc21-aggregated-character-transform}.  The phase
\(\overline{\chi(-h)}\) disappears after taking absolute values.
\end{proof}

In the unmasked case, substituting
\eqref{eq:tpc21-unmasked-B-factorization} into
\eqref{eq:tpc21-masked-fourth-moment} makes it literally a fourth
moment of one-row Dirichlet polynomials, coupled by the shared
\(g\)-character convolution and the outer lcm sum.  In the generic
case it is the square of a masked two-row character sum, hence still a
four-row moment, but not a product of two independent second moments.
The word ``connected'' removes only the principal fiber mode.  The
generic row mask must remain inside
\(\mathcal B_{\mathfrak m}^{u,v}\); why scalar degenerate removal cannot
be performed after this nonlinear norm is explained in
\cref{sec:tpc21-nonlinear}.

\subsection{The Gauss--conductor bridge back to additive frequency}

Put \(e_s(z)=\exp(2\pi iz/s)\), and for
\(\chi\in\mathfrak X(q)\) define the generalized Gauss sum
\begin{equation}
 G_q(\chi;t)
 =\sum_{k\in G_q}\chi(k)e_q(tk).
 \label{eq:tpc21-generalized-Gauss}
\end{equation}
Multiplicative inversion gives the exact bridge
\begin{equation}
 \boxed{
 S_q(r)
 =\frac1{\varphi(q)}
  \sum_{\chi\in\mathfrak X(q)}
  \overline{\chi(-h)}
  \mathcal T_q^{\mathfrak m}(\chi)
  G_q(\chi;r\overline H).}
 \label{eq:tpc21-additive-character-bridge}
\end{equation}

For completeness, the Gauss factor in
\eqref{eq:tpc21-additive-character-bridge} can be evaluated at every
integer, not merely at primitive additive frequencies.

\begin{proposition}[Induced-character Gauss factor]
\label{prop:tpc21-induced-Gauss}
Suppose that \(q\) is squarefree.  Let \(f\mid q\) be the conductor of
\(\chi\), let \(\chi^*\pmod f\) be the inducing primitive character,
and put \(d=q/f\).  Then, for every integer \(t\),
\begin{equation}
 G_q(\chi;t)
 =
 \begin{cases}
  \overline{\chi^*(t)}\chi^*(d)
  \tau_f(\chi^*)c_d(t),&(t,f)=1,\\
  0,&(t,f)>1,
 \end{cases}
 \label{eq:tpc21-induced-Gauss}
\end{equation}
where
\(\tau_f(\chi^*)=\sum_{x\bmod f}\chi^*(x)e_f(x)\), and the
conductor-one conventions are \(\tau_1(\mathbf1)=1\) and
\(\chi^*(d)=1\).
\end{proposition}

\begin{proof}
Factor the sum by CRT modulo \(f\) and \(d\).  The two factors are
\[
 \sum_{x\in G_f}\chi^*(x)e_f(t\overline d\,x)
 \quad\text{and}\quad
 \sum_{y\in G_d}e_d(t\overline f\,y).
\]
The first vanishes when \((t,f)>1\), and otherwise equals
\(\overline{\chi^*(t)}\chi^*(d)\tau_f(\chi^*)\); the second is
\(c_d(t)\).
\end{proof}

As \((H,q)=1\), inserting \eqref{eq:tpc21-induced-Gauss} in
\eqref{eq:tpc21-additive-character-bridge} yields
\begin{align}
 S_q(r)
 ={}&\frac1{\varphi(q)}
 \sum_{\substack{f\mid q\\(f,r)=1}}c_{q/f}(r)
 \sum_{\chi^*\ (\bmod f)}^{*}
 \chi^*(q/f)\tau_f(\chi^*)
 \overline{\chi^*(-hr\overline H)}
 \mathcal T_q^{\mathfrak m}(\chi_q),
 \label{eq:tpc21-conductor-expansion}
\end{align}
where \(\chi_q\) is the character modulo \(q\) induced by \(\chi^*\),
and the star denotes primitive characters.  The \(f=1\) term is exactly
\begin{equation}
 \frac{\mathcal T_q^{\mathfrak m}(\mathbf1_q)}{\varphi(q)}c_q(r)
 =\overline Z_q^{\mathfrak m}c_q(r),
 \label{eq:tpc21-conductor-principal-term}
\end{equation}
so the terms \(f>1\) give the additive discrepancy.

The conductor normalization is important.  If \((r,f)=1\), then
\begin{equation}
 |G_q(\chi;r\overline H)|
 =\sqrt f\,|c_{q/f}(r)|.
 \label{eq:tpc21-Gauss-conductor-size}
\end{equation}
In particular, for \((r,q)=1\) the magnitude is \(\sqrt f\), not
\(\sqrt q\) unless \(\chi\) is primitive modulo \(q\).  Applying the
primitive Gauss magnitude \(\sqrt q\) to every character modulo \(q\)
would therefore be incorrect.

There is also no automatic average gain hidden in the conductors.  A
local count at each prime \(p\mid q\) gives one principal character of
conductor factor \(1\) and \(p-2\) nonprincipal characters of conductor
factor \(p\).  Hence
\begin{equation}
 \sum_{\chi\in\mathfrak X(q)}\operatorname{cond}(\chi)
 =\prod_{p\mid q}\bigl(1+p(p-2)\bigr)
 =\varphi(q)^2.
 \label{eq:tpc21-conductor-mass}
\end{equation}
For \((t,q)=1\), this is equivalently recovered by applying character
Parseval to the unit-modulus function \(k\mapsto e_q(tk)\), which gives
\(\sum_\chi|G_q(\chi;t)|^2=\varphi(q)^2\).  Thus a generic
Cauchy--Schwarz estimate in the character variable merely transfers
the full fiber energy; it does not create cancellation.

\subsection{The naive large-sieve boundary}

The standard multiplicative large sieve states, in one common
normalization, that
\begin{equation}
 \sum_{s\leq S}\frac{s}{\varphi(s)}
 \sum_{\chi\ (\bmod s)}^{*}
 \left|\sum_{n\leq N}a_n\chi(n)\right|^2
 \ll (S^2+N)\sum_{n\leq N}|a_n|^2
 \label{eq:tpc21-standard-character-large-sieve}
\end{equation}
\cite[Chapter~7]{IwaniecKowalski2004}.  This applies directly to one
of the polynomials in \eqref{eq:tpc21-row-Dirichlet-polynomials}.  It
does not, without further work, estimate
\eqref{eq:tpc21-masked-fourth-moment}: that expression contains a
shared-\(g\) convolution, an lcm aggregation, the square of a bilinear
row sum, and the genuinely two-row mask \(\mathfrak m\).  Dropping the
mask to manufacture a product of polynomials changes the nonlinear
quantity, as \cref{sec:tpc21-nonlinear} makes explicit.

At a fixed modulus, complete character orthogonality already gives
\begin{equation}
 \frac1{\varphi(s)}
 \sum_{\rho\in\mathfrak X(s)}
 |\mathcal A_x^{(s)}(\rho)|^2
 =\sum_{c\in G_s}
  \left|\sum_{\substack{\omega\in\mathcal A\\m_\omega\equiv c\ (s)}}
  x_\omega\right|^2.
 \label{eq:tpc21-fixed-character-occupancy}
\end{equation}
Thus the naive second-moment route is exactly a row-occupancy bound
unless additional cancellation in the actual coefficients is proved.
The unconditional occupancy estimate developed in
\cref{sec:tpc21-occupancy-frontier}, when inserted in the inherited
restricted-frequency gate, has relative size
\(X^{o(1)}\sqrt{Y/J}\) on \(q\asymp Y\).  It closes only the
short-joint-period side \(Y<J\) with a fixed margin and supplies no
fixed-power saving in a genuinely long block \(Y\geq JX^\eta\).

Consequently the character change of basis does not by itself produce
a new long-lcm wedge.  Such a wedge would require a theorem for the
actual masked moment \eqref{eq:tpc21-masked-fourth-moment}, or another
arithmetic estimate that retains its row-pair mask and beats the
occupancy scale by a quantified margin.  No such estimate is claimed
here; the large-divisor legs excluded in \eqref{eq:tpc21-theta-choices}
also remain outside this reduction.
